% !TEX program = pdflatex
\documentclass[10pt,aspectratio=169]{beamer}
\usepackage[utf8]{inputenc}
\usepackage[T1]{fontenc}
\usepackage[english]{babel}
\usepackage{lmodern}
\usepackage{microtype}
\usepackage{amsmath,amssymb,bm,mathtools}
\usepackage{booktabs,array}
\usetheme[numbering=fraction,progressbar=frametitle]{metropolis}
\setbeamertemplate{frame numbering}[fraction]
\setbeamertemplate{navigation symbols}{}

\definecolor{CourseBlue}{HTML}{1A3C68}
\definecolor{CourseGreen}{HTML}{216E39}
\definecolor{CourseGold}{HTML}{B8860B}
\definecolor{CourseRed}{HTML}{9A2C2C}
\definecolor{SoftBlue}{HTML}{F2F6FA}
\definecolor{SoftGreen}{HTML}{F1F8F3}
\definecolor{SoftRed}{HTML}{FFF3F3}

\setbeamercolor{normal text}{fg=black,bg=white}
\setbeamercolor{structure}{fg=CourseBlue}
\setbeamercolor{title}{fg=CourseBlue}
\setbeamerfont{title}{size=\Large,series=\bfseries}
\setbeamercolor{frametitle}{bg=CourseBlue,fg=white}
\setbeamerfont{frametitle}{series=\bfseries,size=\large}
\setbeamertemplate{frametitle}{
  \nointerlineskip
  \begin{beamercolorbox}[wd=\paperwidth,ht=3.2ex,dp=1.4ex,leftskip=1em,rightskip=1em]{frametitle}
    \usebeamerfont{frametitle}\insertframetitle
  \end{beamercolorbox}
}
\setbeamercolor{block title}{bg=CourseBlue,fg=white}
\setbeamercolor{block body}{bg=SoftBlue,fg=black}
\setbeamercolor{block title example}{bg=CourseGreen,fg=white}
\setbeamercolor{block body example}{bg=SoftGreen,fg=black}
\setbeamercolor{block title alerted}{bg=CourseRed,fg=white}
\setbeamercolor{block body alerted}{bg=SoftRed,fg=black}
\setbeamertemplate{blocks}[rounded][shadow=false]

\newenvironment{modelbox}{\begin{beamercolorbox}[rounded=true,sep=0.7em,wd=\textwidth]{block body}
  \textbf{\color{CourseBlue}MODEL}\quad}{\end{beamercolorbox}}
\newenvironment{marketbox}{\begin{beamercolorbox}[rounded=true,sep=0.7em,wd=\textwidth]{block body example}
  \textbf{\color{CourseGreen}NO-ARBITRAGE}\quad}{\end{beamercolorbox}}
\newenvironment{simplificationbox}{\begin{beamercolorbox}[rounded=true,sep=0.7em,wd=\textwidth]{block body alerted}
  \textbf{\color{CourseRed}WATCH OUT}\quad}{\end{beamercolorbox}}
\newenvironment{pausebox}{\begin{beamercolorbox}[rounded=true,sep=0.8em,wd=\textwidth]{block body}
  \textbf{\color{CourseGold}YOUR TURN}\quad}{\end{beamercolorbox}}
\newcommand{\E}{\mathbb E}
\newcommand{\Q}{\mathbb Q}
\newcommand{\cF}{\mathcal F}
\newcommand{\dd}{\,\mathrm d}

\AtBeginSection[]{
  \begin{frame}[plain]\vfill
    {\Large\bfseries\color{CourseBlue}\insertsectionhead\par}
    \medskip{\color{CourseBlue}\rule{0.92\textwidth}{0.8pt}}\vfill
  \end{frame}
}

\title{American Options and Optimal Stopping}
\subtitle{When exercising is a decision, not just a payoff}
\author{}
\date{Lecture 10}

\begin{document}
\begin{frame}[plain]
  \vfill
  {\LARGE\bfseries\color{CourseBlue}American Options and Optimal Stopping\par}
  \medskip{\large When exercising is a decision, not just a payoff\par}
  \medskip{\color{CourseGold}\rule{0.94\textwidth}{0.8pt}}
  \medskip{Lecture 10\par}\vfill
\end{frame}

\begin{frame}{One new right changes the pricing problem}
\begin{block}{The question}
An American option lets its holder choose \emph{when} to exercise. How do we
put a no-arbitrage value on a payoff whose exercise date is itself a decision?
\end{block}
\begin{enumerate}
  \item Separate the right to exercise from the payoff received on exercise.
  \item Turn the exercise date into a stopping-time problem.
  \item Solve the discrete problem by backward induction.
  \item Read the continuous-time version as a free-boundary problem.
  \item See how trees, PDEs, and simulation implement the idea.
\end{enumerate}
\end{frame}

\section{The exercise right}

\begin{frame}{European and American options have the same payoff map}
Let \(g(S)=(S-K)^+\) for a call or \(g(S)=(K-S)^+\) for a put.
\begin{columns}[T,onlytextwidth]
\column{0.48\textwidth}
\begin{block}{European option}
Exercise is fixed at maturity:
\[
  \text{payoff}=g(S_T).
\]
There is no exercise decision before \(T\).
\end{block}
\column{0.48\textwidth}
\begin{exampleblock}{American option}
The holder chooses an admissible time
\[
  t\leq\tau\leq T,\qquad \text{payoff}=g(S_\tau).
\]
The payoff map is unchanged; the timing right is valuable.
\end{exampleblock}
\end{columns}
\begin{marketbox}
Maturity is always available. Therefore an American option cannot be worth
less than its European counterpart.
\end{marketbox}
\end{frame}

\begin{frame}{The first inequality is immediate}
For the same strike, underlying, and maturity,
\[
\boxed{V_t^{\mathrm{Am}}\geq V_t^{\mathrm{Eu}}.}
\]
\begin{block}{Why?}
The American holder can choose \(\tau=T\), thereby reproducing the European
payoff. Any better exercise rule can only add value.
\end{block}
\begin{pausebox}
Does \emph{in the money} automatically mean \emph{exercise now} for an
American option?
\end{pausebox}
\pause
\begin{alertblock}{No}
An in-the-money option should be exercised only if \(g(S_t)\) is at least
its \emph{continuation value}.
\end{alertblock}
\end{frame}

\begin{frame}{Early exercise has an economic cost and benefit}
\begin{columns}[T,onlytextwidth]
\column{0.48\textwidth}
\begin{block}{Exercise a put now}
Receive \(K-S_t\) and invest the cash. This can be attractive when the put
is deep in the money and interest rates are positive.
\end{block}
\column{0.48\textwidth}
\begin{exampleblock}{Wait instead}
Keep downside insurance and optionality. A further fall in the stock can
increase the eventual payoff.
\end{exampleblock}
\end{columns}
\begin{simplificationbox}
An American call on a non-dividend-paying stock is a special case: under the
usual Black--Scholes assumptions, early exercise is never optimal. A call on
a dividend-paying stock may be different.
\end{simplificationbox}
\end{frame}

\begin{frame}{A stopping time prevents clairvoyance}
At date \(t\), an exercise rule may use only information known at \(t\).
Formally, \(\tau\) is a stopping time if
\[
 \{\tau=t\}\in\cF_t \qquad \text{for every }t.
\]
\begin{columns}[T,onlytextwidth]
\column{0.48\textwidth}
\begin{exampleblock}{Admissible rule}
Exercise the put the first time the stock is below \(80\).
\end{exampleblock}
\column{0.48\textwidth}
\begin{alertblock}{Not admissible}
Exercise today if tomorrow's stock will be below \(80\).
\end{alertblock}
\end{columns}
\begin{marketbox}
The word \emph{stopping} encodes the information constraint. It makes the
exercise choice economically meaningful.
\end{marketbox}
\end{frame}

\section{Discrete time: value by backward induction}

\begin{frame}{The American price is an optimal-stopping value}
Work on dates \(0,1,\ldots,T\), with one-period gross rate \(R=1+r\) and
discount factor \(\beta=R^{-1}\). Under a risk-neutral measure \(\Q\),
\[
\boxed{V_t=\operatorname*{ess\,sup}_{\tau\in\mathcal T_t}
  \E^\Q\!\left[\beta^{\tau-t}g(S_\tau)\mid\cF_t\right],}
\]
where \(\mathcal T_t\) contains stopping times in \(\{t,\ldots,T\}\).
\begin{block}{Interpretation}
At each node, evaluate every feasible future exercise rule under \(\Q\), then
choose the best one. The rule may adapt to information revealed along the path.
\end{block}
\end{frame}

\begin{frame}{At every node there are exactly two actions}
\begin{columns}[T,onlytextwidth]
\column{0.48\textwidth}
\begin{block}{Exercise}
Stop now and receive intrinsic value
\[
  \boxed{g(S_t).}
\]
No future uncertainty remains.
\end{block}
\column{0.48\textwidth}
\begin{exampleblock}{Continue}
Keep the option for one period:
\[
  \boxed{C_t=\beta\E^\Q[V_{t+1}\mid\cF_t].}
\]
Tomorrow's value already reflects later optimal choices.
\end{exampleblock}
\end{columns}
\begin{marketbox}
Backward induction gives
\[
\boxed{V_T=g(S_T),\qquad V_t=\max\{g(S_t),C_t\}.}
\]
\end{marketbox}
\end{frame}

\begin{frame}{The exercise rule follows the comparison}
\[
\text{exercise at }t
\quad\Longleftrightarrow\quad
g(S_t)\geq \beta\E^\Q[V_{t+1}\mid\cF_t].
\]
\begin{center}
\begin{tabular}{c|c|c}
\toprule
 & Intrinsic \(g(S_t)\) & Continuation \(C_t\)\\
\midrule
Exercise region & \(g(S_t)\geq C_t\) & stop\\
Continuation region & \(g(S_t)<C_t\) & wait\\
\bottomrule
\end{tabular}
\end{center}
\end{frame}

\begin{frame}{Worked example: a two-period American put}
\begin{modelbox}
\begin{center}
\begin{tabular}{c|c}
\toprule
Initial stock & \(S_0=100\)\\
Up/down factors & \(u=1.2,\quad d=0.8\)\\
Gross risk-free return & \(R=1.05\)\\
Strike & \(K=105\)\\
\bottomrule
\end{tabular}
\end{center}
The risk-neutral up probability is
\[
p=\frac{R-d}{u-d}=\frac{1.05-0.8}{1.2-0.8}=0.625,
\qquad 1-p=0.375.
\]
\end{modelbox}
At date \(1\), stock values are \(120,80\). At maturity, the recombining
tree has values \(144,96,64\).
\end{frame}

\begin{frame}{Step 1: write the terminal payoffs}
For the put \(g(S)=(105-S)^+\):
\begin{center}
\begin{tabular}{c|ccc}
\toprule
\(S_2\) & \(144\) & \(96\) & \(64\)\\
\midrule
\(g(S_2)\) & \(0\) & \(9\) & \(41\)\\
\bottomrule
\end{tabular}
\end{center}
\begin{block}{The node with \(S_1=120\)}
Immediate exercise gives \(0\). Continuing gives
\[
C_1^{(120)}=\frac{0.625(0)+0.375(9)}{1.05}=3.21.
\]
So we continue, and \(V_1^{(120)}=3.21\).
\end{block}
\end{frame}

\begin{frame}{Step 2: is early exercise optimal at \(S_1=80\)?}
\begin{pausebox}
At \(S_1=80\), the immediate put payoff is \(105-80=25\). The two terminal
payoffs after waiting are \(9\) and \(41\). Should the holder exercise or
continue?
\end{pausebox}
\pause
\begin{exampleblock}{Compare the two values}
\[
g(80)=25,\qquad
C_1^{(80)}=\frac{0.625(9)+0.375(41)}{1.05}=20.
\]
Since \(25>20\), exercise at this node and set \(V_1^{(80)}=25\).
\end{exampleblock}
\begin{marketbox}
The optimal action depends on the node. The American put has an exercise
region, not a pre-announced exercise date.
\end{marketbox}
\end{frame}

\begin{frame}{Step 3: return to time zero}
At \(t=0\), immediate exercise is worth \(g(100)=5\). Continuing is worth
\[
C_0=\frac{0.625(3.21)+0.375(25)}{1.05}=10.84.
\]
Therefore
\[
\boxed{V_0^{\mathrm{Am}}=\max\{5,10.84\}=10.84.}
\]
\begin{block}{The policy}
Wait at time \(0\); at time \(1\), exercise only in the down node \(S_1=80\).
The price and the policy are computed together.
\end{block}
\end{frame}

\begin{frame}{The tree records both the exercise policy and the hedge}
\begin{center}
\begin{tabular}{c|c|c|c|c}
\toprule
Date & Stock & Exercise value & Continuation value & Decision\\
\midrule
\(0\) & \(100\) & \(5\) & \(10.84\) & continue\\
\(1\), up & \(120\) & \(0\) & \(3.21\) & continue\\
\(1\), down & \(80\) & \(25\) & \(20\) & exercise\\
\bottomrule
\end{tabular}
\end{center}
At time \(0\), the one-period replicating hedge for the two future option
values \(3.21\) and \(25\) uses
\[
\Delta_0=\frac{3.214-25}{120-80}=-0.545,\qquad
B_0\approx65.31.
\]
\begin{marketbox}
\(\Delta_0S_0+B_0=10.84\). The hedge finances the option while it continues;
at the down node it instead finances the exercise payment \(25\).
\end{marketbox}
\end{frame}

\begin{frame}{The early-exercise premium is visible in the numbers}
If exercise were allowed only at maturity, the European put price would be
\[
V_0^{\mathrm{Eu}}
=\frac{2(0.625)(0.375)(9)+(0.375)^2(41)}{1.05^2}
=9.06.
\]
Hence
\[
\boxed{V_0^{\mathrm{Am}}-V_0^{\mathrm{Eu}}=10.84-9.06=1.78.}
\]
\begin{alertblock}{What creates the difference?}
Not a different terminal payoff. It is the right to take \(25\) at the down
node rather than be forced to carry the position to maturity.
\end{alertblock}
\end{frame}

\begin{frame}{A fast decision check}
\begin{pausebox}
At a node of an American put, intrinsic value is \(18\) and discounted
risk-neutral continuation value is \(21\). What is the action and what is the
option value at the node?
\end{pausebox}
\pause
\begin{exampleblock}{Answer}
Continue. The option is worth \(\max\{18,21\}=21\). Being in the money is
not enough: exercising would throw away \(3\) of continuation value.
\end{exampleblock}
\end{frame}

\begin{frame}{The Snell envelope is the discounted value process}
\small
Define discounted exercise payoff and value by
\[
Y_t=\beta^t g(S_t),\qquad \widetilde V_t=\beta^tV_t.
\]
The recursion becomes
\[
\widetilde V_T=Y_T,\qquad
\widetilde V_t=\max\bigl\{Y_t,\E^\Q[\widetilde V_{t+1}\mid\cF_t]\bigr\}.
\]
\begin{block}{Snell envelope}
\(\widetilde V\) is the smallest \(\Q\)-supermartingale that dominates the
discounted payoff process \(Y\).
\end{block}
\begin{marketbox}
It is the American analogue of risk-neutral valuation: the minimal
discounted process that can always cover an exercise request. The next
slide shows this is \emph{both} the buyer's optimal-stopping value and the
seller's minimal hedging cost --- the same number, seen from two sides.
\end{marketbox}
\end{frame}

\begin{frame}{The seller must superhedge every admissible exercise date}
\small
\begin{columns}[T,onlytextwidth]
\column{0.48\textwidth}
\begin{block}{Buyer's value}
The recursion \emph{is} the buyer's best achievable outcome:
\[
V_t=\max\{g(S_t),\beta\E^\Q[V_{t+1}\mid\cF_t]\}.
\]
No stopping rule can do better for the holder.
\end{block}
\column{0.48\textwidth}
\begin{exampleblock}{Seller's price}
The seller does not know which \(\tau\) the buyer will pick. A
self-financing \(X^\pi\) must satisfy
\[
  X^\pi_u\geq g(S_u)
  \qquad\text{for every admissible }u\geq t.
\]
\end{exampleblock}
\end{columns}
The no-arbitrage price is the cheapest such protection:
\[
\boxed{V_t^{\mathrm{Am}}
=\inf\left\{x:\exists\,\pi\text{ self-financing with }
X_u^\pi\geq g(S_u)\text{ for every }u\geq t\right\}.}
\]
\begin{marketbox}
Tracking \(V_t\) with a replicating trade already superhedges every node, so
it costs no more than \(V_t\); no cheaper strategy survives every stopping
choice the buyer might make. Seller's minimal capital = buyer's optimal
value.
\end{marketbox}
\end{frame}

\section{Continuous time: an obstacle and a boundary}

\begin{frame}{In Black--Scholes, the choice is over all future times}
\small
Under \(\Q\), suppose
\[
\dd S_t=rS_t\dd t+\sigma S_t\dd W_t^{\Q}.
\]
\begin{block}{Why work directly under \(\Q\)?}
Black--Scholes pricing uses discounted expectations under the risk-neutral
measure \(\Q\) (Lectures~6--8). The discrete recursion already used this
measure; we now let exercise range over all stopping times.
\end{block}
For an American payoff \(g\), its value function is
\[
\boxed{v(t,s)=\sup_{\tau\in\mathcal T_{t,T}}
\E^\Q\!\left[e^{-r(\tau-t)}g(S_\tau)\mid S_t=s\right].}
\]
\begin{columns}[T,onlytextwidth]
\column{0.48\textwidth}
\begin{block}{Exercise region}
\(v(t,s)=g(s)\): take the payoff now.
\end{block}
\column{0.48\textwidth}
\begin{exampleblock}{Continuation region}
\(v(t,s)>g(s)\): keep the option alive.
\end{exampleblock}
\end{columns}
\end{frame}

\begin{frame}{From an exercise grid to the continuous-time PDE}
First restrict exercise to \(0,h,2h,\ldots,T\), and denote this
\textbf{Bermudan} value by \(v_h\). At a grid date \(t<T\),
\[
\boxed{
v_h(t,s)=\max\left\{g(s),\ \E^\Q_{t,s}
\!\left[e^{-rh}v_h(t+h,S_{t+h})\right]\right\}.}
\]
\begin{block}{Let the exercise mesh tend to zero}
As \(h\downarrow0\), the grid values approach the American value.
In the continuation region, a short-time expansion yields the
Black--Scholes PDE; the stopping branch gives \(v\geq g\).
\end{block}
\begin{simplificationbox}
For a finite \(h\), this recursion excludes exercise inside \((t,t+h)\).
It is not an exact finite-step identity for continuous-time exercise.
\end{simplificationbox}
\end{frame}

\begin{frame}{The PDE has an obstacle}
\begin{columns}[T,onlytextwidth]
\column{0.48\textwidth}
\begin{block}{Exercise branch}
Stopping attains the max: \(v(t,s)=g(s)\). An American value must stay
above immediate exercise, \(v(t,s)\geq g(s)\).
\end{block}
\column{0.48\textwidth}
\begin{exampleblock}{Continuation branch}
Continuing attains the max: expanding
\(\E^\Q[e^{-rh}v(t+h,S_{t+h})]\) to first order in \(h\) via It\^o forces
its \(h\)-coefficient to vanish --- the Black--Scholes PDE.
\end{exampleblock}
\end{columns}
Together these are the variational inequality
\[
\boxed{\max\left\{g(s)-v,\;
v_t+rs\,v_s+\tfrac12\sigma^2s^2v_{ss}-rv\right\}=0.}
\]
The terminal condition is \(v(T,s)=g(s)\).
\begin{simplificationbox}
Unlike a European option, the region where the PDE holds is not known in
advance. Finding it is part of the pricing problem.
\end{simplificationbox}
\end{frame}

\begin{frame}{The put PDE needs terminal and spatial conditions}
Assume no dividends, \(r\geq0\), and \(g(s)=(K-s)^+\). Write
\[
\mathcal Av=rs\,v_s+\tfrac12\sigma^2s^2v_{ss}-rv.
\]
The obstacle problem can be written as three conditions:
\[
\boxed{v-g\geq0,\qquad -(v_t+\mathcal Av)\geq0,\qquad
       (v-g)(v_t+\mathcal Av)=0.}
\]
\begin{block}{Data for the backward calculation}
\[
v(T,s)=(K-s)^+,\qquad v(t,0)=K,\qquad
\lim_{s\to\infty}v(t,s)=0.
\]
On a finite grid, choose \(S_{\max}>K\) and approximate
\(v(t,S_{\max})=0\). Check the effect of increasing \(S_{\max}\).
\end{block}
\textbf{Continue:} \(v>g\) forces the PDE equality.
\quad\textbf{Exercise:} \(v=g\).
\end{frame}

\begin{frame}{For an American put, the boundary is the decision rule}
For a put, there is typically a boundary \(s^*(t)\) such that
\[
\begin{cases}
\text{exercise}, & s\leq s^*(t),\\
\text{continue}, & s>s^*(t).
\end{cases}
\]
At the boundary, the value joins the payoff smoothly:
\[
\boxed{v(t,s^*(t))=K-s^*(t),\qquad v_s(t,s^*(t))=-1.}
\]
\begin{block}{Boundary conditions}
At maturity, \(v(T,s)=(K-s)^+\). For positive rates, the put boundary
typically approaches \(K\) from below as maturity approaches.
\end{block}
\end{frame}

\begin{frame}{Why a non-dividend-paying American call is European}
Let \(C^{\mathrm{Eu}}(t,s)\) be a European call and suppose \(r\geq0\), with
no dividends. For every \(t<T\),
\[
C^{\mathrm{Eu}}(t,s)\geq s-Ke^{-r(T-t)}\geq s-K.
\]
Thus the European call already dominates the value from immediate exercise,
\((s-K)^+\).
\begin{marketbox}
Waiting until maturity is feasible for the American holder, and immediate
exercise never improves on the European value. Therefore
\[
\boxed{C^{\mathrm{Am}}=C^{\mathrm{Eu}}\quad\text{(no dividends).}}
\]
\end{marketbox}
\end{frame}

\begin{frame}{Dividends change the call trade-off}
\begin{columns}[T,onlytextwidth]
\column{0.48\textwidth}
\begin{block}{No dividend}
The option holder does not receive a stock dividend. Giving up the option to
buy the stock early has no compensating cash-flow benefit.
\end{block}
\column{0.48\textwidth}
\begin{exampleblock}{Known dividend}
Exercising before an ex-dividend date can give ownership of the stock and
therefore the dividend. Early exercise of a call can become optimal.
\end{exampleblock}
\end{columns}
\begin{pausebox}
Which object changes with dividends: the payoff, stopping-time logic, or
the continuation value?
\end{pausebox}
\pause
\begin{exampleblock}{Answer}
The payoff and stopping-time logic remain. Dividends alter the economics of
waiting, hence the continuation value and the exercise boundary.
\end{exampleblock}
\end{frame}

\section{Computing American prices}

\begin{frame}{A binomial tree is backward induction made numerical}
With \(N\) steps of length \(\Delta t=T/N\), the CRR parameters are
\[
u=e^{\sigma\sqrt{\Delta t}},\quad d=u^{-1},\quad
p=\frac{e^{r\Delta t}-d}{u-d}.
\]
At every node \((n,j)\), compute
\[
\boxed{V_n^j=\max\left\{g(S_n^j),e^{-r\Delta t}
\left[pV_{n+1}^{j+1}+(1-p)V_{n+1}^j\right]\right\}.}
\]
\begin{block}{Output}
The tree supplies both a price and a discrete approximation to the exercise
boundary: mark every node at which the maximum selects \(g(S_n^j)\).
\end{block}
\end{frame}

\begin{frame}{Finite differences: build the grid and start at maturity}
For the no-dividend put, with \(r\geq0\), set
\[
s_i=i\Delta s,\quad \Delta s=S_{\max}/M,\qquad
t_n=n\Delta t,\quad \Delta t=T/N,\qquad V_i^n\approx v(t_n,s_i).
\]
Start with \(V_i^N=g_i=(K-s_i)^+\) and work
\textbf{backward}: \(n=N-1,\ldots,0\).
\begin{block}{Centered differences at an interior stock node}
\[
v_s(t_n,s_i)\approx\frac{V_{i+1}^n-V_{i-1}^n}{2\Delta s},\qquad
v_{ss}(t_n,s_i)\approx
\frac{V_{i+1}^n-2V_i^n+V_{i-1}^n}{(\Delta s)^2}.
\]
\end{block}
\begin{block}{Boundary values and the obstacle}
At every time level impose \(V_0^n=K\), \(V_M^n=0\).
The obstacle is \(V_i^n\geq g_i\); the price is read at \(t_0=0\).
\end{block}
\end{frame}

\begin{frame}{Explicit scheme: a continuation step, then the maximum}
\small
Discretizing \(v_t+\mathcal Av=0\) with spatial terms at the known
level \(n+1\) gives, for \(i=1,\ldots,M-1\),
\[
\boxed{C_i^n=a_iV_{i-1}^{n+1}+b_iV_i^{n+1}+c_iV_{i+1}^{n+1},
\qquad V_i^n=\max\{g_i,C_i^n\}.}
\]
Here \(s_i=i\Delta s\), so
\[
\begin{aligned}
a_i&=\tfrac{\Delta t}{2}(\sigma^2i^2-ri),&
c_i&=\tfrac{\Delta t}{2}(\sigma^2i^2+ri),\\
b_i&=1-\Delta t(\sigma^2i^2+r),&
a_i+b_i+c_i&=1-r\Delta t.
\end{aligned}
\]
\begin{alertblock}{Check the weights before using this explicit scheme}
Require \(a_i,b_i,c_i\geq0\) at every interior node. For this centered stencil:
\[
\sigma^2i^2\geq ri,\qquad
\Delta t\leq\frac{1}{\sigma^2(M-1)^2+r}.
\]
If a spatial weight is negative, change the drift stencil; reducing
\(\Delta t\) alone cannot fix its sign.
\end{alertblock}
\end{frame}

\begin{frame}{One time step: exercise at 80, continue at 100}
Take \(K=100\), \(r=5\%\), \(\sigma=30\%\),
\(\Delta s=5\), \(\Delta t=1/252\), and \(S_{\max}=200\).
At the last step, the known values are \(V_i^N=(100-s_i)^+\).
\medskip
\begin{center}
\begin{tabular}{c|c|c|c|l}
\toprule
\(s_i\) & \((V_{i-1}^N,V_i^N,V_{i+1}^N)\) & \(C_i^{N-1}\) & \(g_i\) & Decision\\
\midrule
\(80\) & \((25,20,15)\) & \(19.9802\) & \(20\) & Exercise\\
\(100\) & \((5,0,0)\) & \(0.3472\) & \(0\) & Continue\\
\bottomrule
\end{tabular}
\end{center}
\begin{block}{The PDE step and the stopping decision are both visible}
At \(s=80\), \(C=20-rK\Delta t<20\): the maximum restores \(20\).
At \(s=100\), \(C=5a_{20}>0\): keeping the option has value.
These are grid approximations one time step before maturity.
\end{block}
Mark nodes where \(V_i^n=g_i>0\) to approximate the put exercise region.
\end{frame}

\begin{frame}{Implicit scheme: solve the PDE and obstacle together}
\small
Using spatial terms at the unknown level \(n\), define the residual
\[
R_i^n=d_iV_i^n-a_iV_{i-1}^n-c_iV_{i+1}^n-V_i^{n+1},
\qquad d_i=1+\Delta t(\sigma^2i^2+r).
\]
With the same boundary values and \(a_i,c_i\) as before, solve
\[
\boxed{V_i^n-g_i\geq0,\quad R_i^n\geq0,\quad
       (V_i^n-g_i)R_i^n=0.}
\]
\begin{block}{A discrete complementarity problem}
Use projected Gauss--Seidel, projected SOR, or policy iteration until
the constraints hold. With \(a_i,c_i\geq0\), the implicit scheme avoids
the explicit time-step restriction; accuracy still requires refinement.
\end{block}
\begin{simplificationbox}
One unconstrained implicit solve followed by a single maximum is generally
only a splitting approximation, not the solution of this coupled system.
\end{simplificationbox}
\end{frame}

\begin{frame}{Why plain Monte Carlo is not enough}
For a European option, simulation estimates
\[
e^{-rT}\E^\Q[g(S_T)]
\]
by averaging terminal discounted payoffs.
\begin{alertblock}{The American difficulty}
At an intermediate state, the relevant object is
\[
\E^\Q[\text{future cash flow}\mid S_t=s],
\]
the conditional continuation value. Simulated paths do not reveal this
function directly at a single node.
\end{alertblock}
\begin{marketbox}
American pricing is not just simulation; it is simulation plus a way to
estimate the decision rule.
\end{marketbox}
\end{frame}

\begin{frame}{Longstaff--Schwartz learns the continuation value}
Least-squares Monte Carlo works backward through simulated paths.
\begin{enumerate}
  \item Simulate risk-neutral stock paths and start from terminal payoff.
  \item At each earlier date, retain in-the-money paths.
  \item Regress discounted future cash flows on basis functions, e.g.
  \(1,S_t,S_t^2\).
  \item Compare estimated continuation with immediate exercise and update
  the cash-flow date path by path.
  \item Discount the resulting cash flows and average them.
\end{enumerate}
\begin{block}{Interpretation}
Regression does not price the option directly; it estimates the continuation
rule needed to decide whether each path should stop.
\end{block}
\end{frame}

\begin{frame}{Method choice follows the state space}
\begin{center}
\begin{tabular}{p{0.23\textwidth}|p{0.29\textwidth}|p{0.30\textwidth}}
\toprule
Method & Best fit & Main caution\\
\midrule
Binomial tree & One underlying; teaching and robust benchmarks & Many steps may be needed for convergence\\
Finite differences & Low-dimensional Markov models; boundary analysis & Grid size grows quickly with dimension\\
Longstaff--Schwartz & Several risk factors or path simulation available & Quality depends on paths, basis, and regression design\\
\bottomrule
\end{tabular}
\end{center}
\begin{marketbox}
All three methods implement the same comparison: exercise value versus
risk-neutral continuation value.
\end{marketbox}
\end{frame}

\begin{frame}{Three checks catch many implementation mistakes}
Before trusting a computed American price, test the output against consequences
of the theory.
\begin{center}
\begin{tabular}{p{0.28\textwidth}|p{0.56\textwidth}}
\toprule
Check & What should be true\\
\midrule
Exercise right & \(V^{\mathrm{Am}}\geq V^{\mathrm{Eu}}\) for the same contract.\\
No-dividend call & \(C^{\mathrm{Am}}=C^{\mathrm{Eu}}\), up to numerical tolerance.\\
Refinement & A tree or grid should stabilize as \(N\) increases or the mesh is refined.\\
\bottomrule
\end{tabular}
\end{center}
\begin{pausebox}
If numerical output gives \(V^{\mathrm{Am}}<V^{\mathrm{Eu}}\), is that an
economic result or a bug?
\end{pausebox}
\pause
\begin{alertblock}{Debug it}
Exercise at maturity is feasible, so an American value cannot be lower.
\end{alertblock}
\end{frame}

\section{Takeaways}

\begin{frame}{The whole lecture in one recursion}
\[
\boxed{V_t=\max\left\{\underbrace{g(S_t)}_{\text{exercise}},
\underbrace{\beta\E^\Q[V_{t+1}\mid\cF_t]}_{\text{continue}}\right\}.}
\]
\begin{enumerate}
  \item An American option embeds a stopping right, so its value is at least
  the corresponding European value.
  \item The decision is dynamic: optimal exercise is a region of states.
  \item In continuous time, that region becomes a free boundary and the PDE
  becomes an obstacle problem.
  \item Trees, PDE solvers, and LSMC differ computationally, not conceptually.
\end{enumerate}
\end{frame}

\begin{frame}{Final check: read the economics before the formula}
\begin{pausebox}
You are handed an American put and told only that it is deep in the money,
rates are positive, and little time remains. Which comparison should you make
before declaring that exercise is optimal?
\end{pausebox}
\pause
\begin{exampleblock}{Answer}
Compare intrinsic value \(K-S_t\) with discounted risk-neutral continuation
value. The clues make early exercise plausible, but the Bellman comparison
makes the decision.
\end{exampleblock}
\vfill
\begin{center}
{\large\color{CourseBlue}\bfseries Next: using these tools to value and hedge exercise features.}
\end{center}
\end{frame}

\end{document}
